% English manuscript fragment. Source: 01_英文正文.md.
% Historical v1 fragment only; the current manuscript is in ../重构稿_v2/.
% The separate bibliography database was removed at the user's request on 2026-09-25.
% Reference information remains in 01_英文正文.md; this fragment is not a standalone build target.
% No document class, author metadata, or bibliography style is imposed here.

% Suggested title: DiGO: Diffusion-Driven Graph Optimization for Bundle Recommendation
\begin{abstract}
Bundle recommendation seeks to match users with sets of related items by combining bundle-level interactions with item-level information. Despite progress in multi-view modeling, learning bundle preferences remains challenging when observed user-bundle (U-B) interactions contain noise and individual users provide limited bundle feedback. We propose DiGO, a diffusion-driven graph optimization framework that couples U-B graph denoising with cross-view preference learning. Its core is a unified pipeline of pretraining, user-conditioned latent diffusion, and U-B graph reconstruction. We first pretrain user and bundle representations on observed U-B interactions and aggregate each user's historical bundle representations into a latent preference vector. A diffusion model learns to recover this vector from progressively perturbed states, conditioned on the pretrained user representation. An auxiliary bundle-set consistency loss keeps the denoising process connected to the user's historical bundle set. The refined preference vector then re-scores observed interactions to select edges for reconstructing the U-B bipartite graph. To address limited bundle feedback, we further introduce pre-fusion cross-view contrastive learning anchored in representations learned from the reconstructed U-B graph. User-side and bundle-side objectives connect this view to the user-item and bundle-item views before their representations are fused for recommendation. Experiments on three public benchmarks show that DiGO outperforms the compared baselines. Component and interaction-selection analyses support the proposed designs, while the user-group analysis shows that the cross-view objective yields the largest relative gains for users with the fewest bundle interactions.

\end{abstract}

\section{Introduction}

Users often seek a coherent set of items rather than an isolated product: a coordinated outfit or a themed playlist can satisfy a need that individual recommendations leave users to assemble themselves. Bundle recommendation addresses this demand by recommending related items as a unit \cite{ma2022crosscbr,sun2024intent}. Its central task is not simply to identify appealing items, but to determine which bundles match a user's preferences. In this work, we focus on ranking predefined bundles using historical interactions and bundle composition information.

This setting involves three relations with distinct roles. User-bundle (U-B) interactions record users' choices of complete bundles and directly support learning the target preference. User-item (U-I) interactions provide evidence of preferences for individual items, while bundle-item (B-I) affiliations specify how items constitute bundles. Representations derived from these relations offer different views of users and bundles: U-B reflects bundle-level collaborative behavior, whereas U-I and B-I provide complementary behavioral and compositional information. The resulting challenge is to learn from these relations in a way that serves bundle ranking, rather than treating them as interchangeable sources of evidence \cite{chen2019dam,chang2020bgcn,ma2024multicbr}.

Research has progressively developed mechanisms to exploit these relations. Attention-based approaches aggregate constituent items and jointly learn from bundle and item interactions \cite{chen2019dam}. Graph neural networks (GNNs) extend this modeling to higher-order connectivity among users, bundles, and items \cite{chang2020bgcn}. More recently, contrastive learning (CL) has enabled cooperation between views: CrossCBR links bundle-level and item-level representations through cross-view objectives \cite{ma2022crosscbr}, while MultiCBR separately models all three relations, fuses their representations, and applies contrastive learning after fusion \cite{ma2024multicbr}. These advances demonstrate the value of auxiliary information, but do not remove the need to examine the U-B evidence on which bundle preference learning relies. We focus on two challenges concerning the quality and availability of this evidence.

\textbf{Noise in user-bundle interactions.} Observed implicit feedback does not always faithfully reflect user preferences \cite{wang2021adt}. In bundle recommendation, this matters because U-B edges both record choices of complete bundles and connect users to the bundle neighborhoods used for graph propagation. Retaining every observed connection without assessing its preference relevance can allow noisy interactions to influence the learned representations, a vulnerability also identified in graph-based recommendation \cite{wu2021sgl}. However, discarding edges alone is insufficient: a user's interacted bundles collectively contain useful preference information. The challenge is to use this historical context to reassess individual interactions, suppressing noise while preserving the preference information needed for bundle ranking.

\textbf{Sparse user-bundle feedback.} For users with few observed bundle interactions, the available history provides limited evidence for learning personalized bundle preferences. U-I behavior and B-I composition can help, but their relation to bundle choice must be learned: liking individual items is not equivalent to preferring a complete bundle, and a bundle's composition alone does not establish a user's interest in it. Prior cross-view learning and distillation methods demonstrate the importance of connecting these information sources \cite{ma2022crosscbr,ren2023dgmae}. This motivates additional learning constraints that relate auxiliary-view representations to bundle-level preferences, so that users with limited U-B feedback can benefit from the available item-level information.

To address these challenges, we propose \textbf{DiGO}, a diffusion-driven graph optimization framework for bundle recommendation. The central design connects \textbf{pretraining, latent diffusion denoising, and U-B graph reconstruction} in one pipeline. We first pretrain user and bundle embeddings on observed U-B interactions, then pool each user's historical bundle embeddings to obtain a latent bundle-preference vector. Drawing on the denoising formulation of diffusion models \cite{ho2020ddpm,wang2023diffrec}, we progressively perturb this vector and train a denoiser to recover it using the pretrained user embedding as a personalized condition. An auxiliary bundle-set consistency loss (SetLoss) encourages the recovered vector to remain consistent with the user's historical bundle representations. The resulting refined preference vector is used to re-score observed U-B interactions and retain high-scoring connections. In this way, preference denoising is translated into a reconstructed U-B bipartite graph for downstream recommendation, rather than ending at the representation level.

On top of this graph, DiGO introduces \textbf{U-B-anchored pre-fusion cross-view contrastive learning} to strengthen preference learning under limited feedback. The recommendation backbone learns user and bundle representations from the reconstructed U-B graph and the U-I and B-I relations, then fuses the view-specific representations for prediction. Before fusion, we use the reconstructed U-B view as the common anchor and construct U-B--U-I and U-B--B-I contrastive pairs for both users and bundles. These objectives encourage agreement between representations of the same entity across the paired views while distinguishing different entities. They connect auxiliary information to the refined bundle-level view through additional learning signals derived from existing data. Thus, graph reconstruction and cross-view learning play complementary roles: the former reassesses the U-B connections used for propagation, and the latter helps the model learn beyond each user's limited bundle history.

Our contributions are summarized as follows:

\begin{itemize}
\item We propose a diffusion-driven U-B graph denoising framework that connects pretrained collaborative representations, user-conditioned latent denoising, and observed-edge selection. A bundle-set consistency constraint supports the denoising process by retaining its connection to historical bundle preferences.
\item We introduce pre-fusion cross-view contrastive learning anchored in the reconstructed U-B view, linking bundle-level representations to U-I and B-I representations on both the user and bundle sides.
\item We evaluate DiGO on three public benchmarks. Overall comparisons, component ablations, and interaction-selection analyses support its effectiveness; a user-group analysis shows the largest relative gains from the proposed cross-view objective in the lowest-interaction group.
\end{itemize}

\section{Related Work}

\subsection{Bundle Recommendation}

Bundle recommendation includes both ranking existing bundles and generating new item combinations \cite{sun2024survey}. For the former, which is our setting, a key modeling issue is how to combine evidence about complete bundles with information about their constituent items. DAM addresses this issue through attentive item aggregation and joint learning of U-B and U-I interactions \cite{chen2019dam}. BGCN extends representation learning to graph neighborhoods, combining bundle-level collaborative information with item-level information \cite{chang2020bgcn}. The broader graph-based recommendation literature, exemplified by LightGCN, establishes neighborhood aggregation as a useful mechanism for learning collaborative representations \cite{he2020lightgcn}. Subsequent bundle models enrich what these representations capture: MIDGN disentangles user intents across global and local views \cite{zhao2022midgn}, while BundleGT uses a hierarchical graph transformer to model bundling strategies \cite{wei2023bundlegt}. These approaches illustrate why a bundle cannot be adequately understood as an item identifier alone.

Another line of work emphasizes cooperation across information sources. CrossCBR contrasts representations from bundle and item views to encourage mutual enhancement \cite{ma2022crosscbr}. MultiCBR explicitly models the B-I graph alongside U-B and U-I graphs, and adopts representation fusion followed by contrastive learning to capture multi-view preference information \cite{ma2024multicbr}. DGMAE takes a different route, distilling first- and higher-order U-B knowledge into an item-view graph masked autoencoder \cite{ren2023dgmae}. Together, these studies show that both the construction of individual views and the mechanism connecting them matter. Building on this multi-view setting, DiGO concentrates on refining the U-B graph through a complete pretraining--diffusion--reconstruction pipeline and using its resulting representations to organize cross-view learning. This focus complements richer intent or composition modeling without requiring a change to the predefined-bundle ranking task.

\subsection{Denoising and Diffusion for Recommendation}

Denoising methods address the mismatch between observed interactions and the preferences a recommender aims to learn. Adaptive Denoising Training uses truncated or reweighted losses to reduce the influence of potentially noisy implicit feedback \cite{wang2021adt}. This line of work motivates assessing interaction relevance rather than treating every observed event as equally informative. Diffusion models offer a different modeling tool: learning to reverse a sequence of controlled corruptions \cite{ho2020ddpm}. In recommendation, DiffRec reconstructs interaction vectors through denoising and also introduces a latent-space variant for efficiency \cite{wang2023diffrec}. DiffuRec instead operates on target-item embeddings, using interaction sequences to condition their recovery for next-item prediction \cite{li2023diffurec}. The choice of denoising object therefore determines how diffusion contributes to the recommendation task.

Diffusion has also been explored for bundle recommendation. DCBR combines conditional U-B interaction denoising with disentangled contrastive learning \cite{li2025dcbr}. RDiffBR applies residual diffusion to item-level bundle embeddings to accommodate changes in bundle composition \cite{zhang2026rdiffbr}. DiGO models an aggregate of historical bundle embeddings conditioned on a pretrained user embedding, and uses the refined vector to select connections from the observed U-B graph. Its emphasis is the connection from personalized latent denoising to the graph used for recommendation: bundle-set consistency regularizes preference recovery, while edge selection translates the recovered preferences into graph structure.

\subsection{Contrastive Learning for Recommendation}

Contrastive learning constructs learning signals by increasing agreement between related representations while distinguishing unrelated ones. SimCLR demonstrates the importance of view construction and the contrastive objective in representation learning \cite{chen2020simclr}. For graph-based recommendation, SGL generates graph views through perturbations such as edge or node dropout and contrasts representations of the same node across views \cite{wu2021sgl}. In knowledge-aware recommendation, KGCL uses consistency under knowledge-graph perturbations to guide interaction augmentation \cite{yang2022kgcl}, while MCCLK connects collaborative, semantic, and structural representations through multi-level cross-view objectives \cite{zou2022mcclk}. These methods make view selection and cross-view correspondence explicit parts of representation learning.

In bundle recommendation, the U-B, U-I, and B-I relations provide task-relevant views without requiring external knowledge graphs. The bundle-specific methods reviewed above differ in how they connect these views and where contrastive learning is applied \cite{ma2022crosscbr,ma2024multicbr}. DiGO builds on this principle with a focused pre-fusion objective: representations learned from the reconstructed U-B graph are paired with the corresponding U-I and B-I representations on both entity sides. The reconstructed U-B view provides a shared reference for jointly optimizing the view-specific representations. This design connects the graph denoising module to auxiliary-view learning while leaving representation fusion to the recommendation backbone.
